% period_ladder.tex
% GeoVac synthesis: the period ladder
% Status: front door for readers of the periods / Feynman-integral literature
% Audience: number theorists and Feynman-integral practitioners; readable by quantum chemists
% This document makes no claim of its own: every statement restates a claim
% owned by a paper of the series, with the owner cited and its tier tag kept.

\documentclass[aps,pra,twocolumn,superscriptaddress,longbibliography]{revtex4-2}

\usepackage{amsmath,amssymb,amsthm,graphicx,xcolor,bm,braket,booktabs}
\usepackage{hyperref}

% Claim-registry and numeric-registry citation macros (as in the other
% syntheses): each renders its literal only; the key is a foreign key.
\newcommand{\gvc}[2]{#2}
\newcommand{\gvq}[2]{#2}
\newcommand{\mHa}{\,\mathrm{mHa}}

\begin{document}

\title{The Period Ladder:\\
The Integrals of Quantum Chemistry Sorted by the Kind of Number They Make}
\thanks{A synthesis in the Geometric Vacuum series}

\author{J.~Loutey}
\affiliation{Independent Researcher, Kent, Washington}

\date{November 19, 2026}

%% ========================================================================
\begin{abstract}
The integrals of molecular quantum chemistry over exponential
(Slater-type) orbitals can be sorted by the class of number they produce.
This note tells that sorting end to end, from the hydrogen atom's sphere
to the banana integrals of quantum field theory, for a reader who works
on periods and Feynman integrals;\ a quantum chemist should be able to
follow it too.  It makes no new claim.  Every statement restates a result
owned by a paper of the Geometric Vacuum series, with the owner cited and
its tier tag kept.  The story has four steps.  The hydrogen atom, by
Fock's projection, is a compact problem whose labels and degeneracies are
integers.  Paper~18's claim is that each released axis carries an
identifiable transcendental price.  Releasing axes by geometry alone (a
second nucleus, a third, a fourth, bending, an external field, a change
of shape) reaches an elliptic period and, across the doors the atlas
records, stops there (an observation, not a theorem).  The class climbs
to a higher rung (rung $n$ is the class of the operator of a Bessel
moment $F_n$) only when one integral links more electron densities:\ a tree of $n$ linked two-center densities has a
holomorphic period whose one-scale operator, for $n=2$ to $5$, is the
Hadamard product (the term-by-term product of the coefficients of their
power-series solutions) of the Legendre operator with the operator of the
equal-mass banana integral with $n-1$ propagators.  Most of the
mathematics on the ladder is prior art, credited in
Sec.~\ref{sec:prior};\ what the series adds is the route from molecular
integrals to it, listed there with its caps.  The note ends with the open
questions and with where each claim can be checked.
\end{abstract}

\maketitle

%% ========================================================================
\section{The skeleton:\ hydrogen on a sphere}
\label{sec:skeleton}
%% ========================================================================

\textbf{[PRIOR ART]} Fock showed in 1935 that stereographic projection
turns the momentum-space hydrogen Schr\"odinger equation into a
rotation-invariant integral equation on the unit three-sphere
$S^3$~\cite{fock1935}.  Its kernel is diagonal on the hyperspherical
harmonics, the eigenfunctions of the Laplace--Beltrami operator
$\Delta_{S^3}$ with eigenvalue $-(n^2-1)$, and the bound-state
Hamiltonian is a function of that operator,
$H=-\tfrac{Z^2}{2}(1-\Delta_{S^3})^{-1}$, so that $E_n=-Z^2/(2n^2)$
(Paper~7~\cite{loutey_paper7}).  The sphere is compact, so everything
that labels a state is a whole number:\ the labels $(n,l,m)$, the
degeneracy $n^2$ per shell, and the Casimir $n^2-1$ of the symmetry group
$SO(4)$.  This label set, with the Casimir as a weight on its nodes, is
what the series calls the skeleton.  Paper~0 reaches the angular part of
the same label count from a packing construction and its structural
choices (motivated, not forced) \textbf{[SYMBOLIC PROOF]};\ the grouping of shells into the principal quantum number is an
\textbf{[OBSERVATION]} there, not forced by its
axioms~\cite{loutey_paper0}.  The matrix entries, degeneracies and
coupling coefficients on the skeleton are integers or rationals;\ no
$\pi$ appears until something is projected off it
(Paper~18~\cite{loutey_paper18}).

One wall must be stated plainly, because an earlier reading of the
series said otherwise.  \textbf{[WALL]} \gvc{graph-s3-wall}{The edge
Laplacian $L=D-A$ of the label lattice does \emph{not} converge to the
Laplace--Beltrami operator of $S^3$:\ its edges never change $l$, so it
splits into $n_{\max}$ disconnected $l$-blocks, and its kernel is
$n_{\max}$-dimensional at every cutoff while $\Delta_{S^3}$ has a single
zero mode} (Paper~7~\cite{loutey_paper7},
\texttt{tests/test\_graph\_s3\_nonconvergence.py}).  What is exact at
every cutoff is the label set and the diagonal Casimir.  The equivalence
with the Schr\"odinger problem is carried by Fock's continuum projection,
not by a graph limit.

The energies themselves are not on the sphere.  \textbf{[SYMBOLIC
PROOF]} The levels $E_n$ enter only through the energy-shell constraint
$p_0^2=-2E$, which sets the scale of the stereographic map
(Paper~7~\cite{loutey_paper7}).  That scale is the mass on every orbital
line of the graphs below.

%% ========================================================================
\section{Paper~18's claim:\ each released axis carries a price}
\label{sec:rule}
%% ========================================================================

A compact group has a discrete spectrum (the Peter--Weyl
theorem~\cite{peter_weyl1927}).  The reading that compactness is the
\emph{source} of spectral discreteness is not original to this series,
and Paper~18 does not claim it:\ it is stated causally for the gauge
group of loop quantum gravity by Liu and Noui~\cite{liu_noui2017}, and earlier and more generally in Rovelli's relational quantum
mechanics~\cite{rovelli1996}.  Nor is the converse true.  Maz'ya
and Shubin give conditions under which a Schr\"odinger operator on
non-compact $\mathbb{R}^n$ has purely discrete
spectrum~\cite{mazya_shubin2005};\ a confining potential suffices, and
the harmonic oscillator is the standing counterexample.  The correct
general statement is at the level of the operator:\ discrete spectrum
with finite multiplicities is equivalent to compactness of the resolvent
(Paper~18~\cite{loutey_paper18}).

What Paper~18 adds is not the mechanism but the \emph{accounting}:\ the
claim that each released axis carries an identifiable transcendental
price, and that those prices can be classified.  Transcendental numbers
enter when the discrete description is projected onto a non-compact
coordinate system, and the paper sorts them by the kind of mismatch the
projection mediates.  Its strongest form, that the transcendental content
of any observable is determined entirely by the type of projection, is
stated there as a falsifiable claim that is under strain from one
observable and is not reconciled with the paper's weaker, route-indexed
form.  We use only the accounting.

Two terms will recur.  A \emph{period}, in the sense of Kontsevich and
Zagier~\cite{kontsevich_zagier2001}, is a number whose real and imaginary
parts are integrals of algebraic functions over domains defined by
polynomial inequalities with rational coefficients.  The \emph{weight} of
a period of polylogarithm type counts its iterated integrations:\ a
logarithm has weight one, a dilogarithm or $\pi^2$ weight two.  The
chemistry integrals below also carry exponentials of the internuclear
distance and the exponential integral $E_1$, so ``number class'' is used
in the wider sense of Paper~18:\ the functions and constants a closed
form needs.

\emph{The sphere's own numbers.}  Before any nucleus is added, an
observer reading the sphere already pays.  Paper~55 takes a census of the
spectral-zeta and heat-trace data of $S^3$ and proves, over that census,
that the transcendentals it identifies sit in three rings:\ $\mathbb{Q}[\pi,\pi^{-1}]$ for the volume slot,
$\bigoplus_k\pi^{2k}\,\mathbb{Q}$ for the heat-kernel coefficients, and
cyclotomic mixed-Tate periods at level at most four, with Catalan's
constant $G=\beta(2)$, for parity-restricted Dirac sums
(Papers~55 and~28~\cite{loutey_paper55,loutey_paper28}).  The indexing of
these rings by the slots of one Mellin transform is a reading, not a
classification, and the paper says that the chemistry-side numbers and
the elliptic period of Sec.~\ref{sec:three} lie outside its census.
Those are the subject of the rest of this note.

%% ========================================================================
\section{The doors:\ what geometry alone can buy}
\label{sec:doors}
%% ========================================================================

The series keeps a register, the atlas (\texttt{docs/atlas.md}), with one
row for each way out of the compact skeleton:\ what the door releases,
what survives, the number class that enters, and a verdict.  The
verdicts are catalogue labels:\ \emph{forced} (derived, with a test that
would fail if it were false), \emph{free} (an input consumed from
outside, with its number class named), \emph{wall} (a proven
obstruction), and \emph{observation} (a measured fact with no
derivation).  They are bookkeeping, not statements about what is
fundamental.  A \emph{rung} is a class of period named by a differential
operator:\ rung $n$ is the class of the operator of the Bessel moment
$F_n$ of Sec.~\ref{sec:climb}, and rung two is the elliptic one.  Table~\ref{tab:doors} selects the doors that change the
geometry, with the atlas numbers so that a reader can find the full row.

\begin{table*}[t]
\caption{The geometric doors, selected from the atlas (which has 32
rows).  ``Price'' is the number class that enters.  Rows 28 to 30 are
recorded in the atlas and its test only, and almost all of their content
is classical prior art credited in the atlas rows.  The lower block lists
the doors that link densities (Sec.~\ref{sec:climb}).}
\label{tab:doors}
\small
\begin{tabular}{@{}rp{3.6cm}p{4.3cm}p{5.3cm}p{2.9cm}@{}}
\toprule
\# & Door & What is released & Price (number class) & Verdict \\
\midrule
1 & the label set $(n,l,m)$ with the Casimir $n^2-1$ &
  nothing:\ this is the skeleton &
  zero:\ integers and rationals & forced \\
2 & the edge Laplacian as a route to $\Delta_{S^3}$ &
  a continuum limit was attempted &
  algebraic numbers, and no limit & wall \\
11 & a second nucleus &
  $l$ across centers;\ $m$ survives &
  weight one per Neumann term, $\pi$-free, no dilogarithm;\ the
  unequal-exponent exchange sum is weight two for $0<a<b$ (Paper~58) & forced (the $m$ rule) \\
13 & a third nucleus with a shared center, $(XY|XZ)$ &
  the third focus &
  an elliptic period $K(m)$ for the radial core &
  wall:\ no closed form of the radial core through factorization, and the
  elliptic-dilogarithm route does not engage;\ a closed form is open \\
14 & a fourth distinct center, $(XY|ZW)$ &
  a fourth center &
  no genus:\ the same genus-one family &
  forced (pre-registered;\ measured, nearly a corollary) \\
15 & bending a triatomic &
  $m$:\ a pair $|m|=\Lambda$ opens at order $s^{2\Lambda}$ &
  exact rationals for the order and the angular factors;\ the slopes are
  exponentials in $R$, no $\pi$ & forced (order, labels);\ slopes measured \\
28 & a uniform electric field on hydrogen &
  $n$ and $l$:\ levels become resonances;\ $m$ survives &
  exact rationals order by order, in a divergent series;\ a trans-series
  price, not a period & free (the field);\ forced (the coefficients) \\
29 & a uniform magnetic field on hydrogen &
  the $n$-degeneracy;\ $l$ mixes;\ $m$ and $z$-parity survive &
  exact rationals order by order, Borel summable;\ the class of an
  eigenvalue at finite field is open & free (the field);\ high field open \\
30 & a quotient of the sphere (lens spaces) &
  nothing physical:\ a symmetry sector is kept &
  integer multiplicities with a rational generating function;\ spectral
  sums in $\bar{\mathbb{Q}}[\pi]$ & forced as mathematics;\ not a physical door \\
31 & curving a tube of hydrogen rings into a torus &
  the tube winding $m$, leaking as $(r/R)^{|\Delta m|}$ &
  none measured:\ the leak order is an integer & observation \\
32 & squashing a hydrogen ring into an ellipse &
  the ring momentum $k$, at order $\epsilon^{p(c)}$ &
  none measured:\ the orders are integers & observation \\
\midrule
17 & a three-electron chain &
  a second momentum integration &
  a K3 surface;\ the holomorphic period is a sum of products of elliptic
  periods & forced;\ generic content observation \\
20 & four linked densities &
  a third momentum integration &
  on the one-scale slice, a fourth-order Calabi--Yau-type operator (AESZ
  \#58), not built from elliptic periods & forced \\
21 & a Coulomb ring of three densities &
  the loop momentum, which integrates out &
  the elliptic rung, below the chain;\ at one center, dilogarithms &
  forced (Coulomb kernels) \\
\bottomrule
\end{tabular}
\end{table*}

Read the upper block of the table from top to bottom.  One center costs
nothing.  A second nucleus releases $l$ and keeps $m$, and its price is
of weight one term by term (Sec.~\ref{sec:two});\ that describes the
integrated value, while the two-center exchange fibre is already
elliptic.  A third nucleus brings an elliptic period
(Sec.~\ref{sec:three}), and a fourth adds no genus:\
\textbf{[MEASURED]} the general four-center integral $(XY|ZW)$ over $1s$
Slater orbitals reduces to one radial integral whose integrand carries
exactly the two dispersion roots of the three-center curve, so a fourth
distinct center adds no genus (Paper~63~\cite{loutey_paper63}).  Bending
releases $m$ at a counted order, \textbf{[SYMBOLIC + MEASURED]} with
exact rational angular prices $6/5$ ($\Lambda=1$) and $5/6$
($\Lambda=2$) on a single shell and a free coefficient whose radial
weights carry exponentials in $R$ and no $\pi$;\ this is the Renner
effect~\cite{renner1934} (Paper~32~\cite{loutey_paper32}).  A uniform
field is an input from outside the skeleton (Paper~34 classes the Stark
and Zeeman fields as such inputs~\cite{loutey_paper34}), and its
perturbation coefficients are rational.  A free quotient of the sphere
is mathematics, not a physical door.  Curving or squashing a ring of
hydrogen atoms breaks a winding label as a power of the deformation, and
the power is an integer \textbf{[MEASURED, pre-registered;\ GO]}
(Paper~62~\cite{loutey_paper62};\ the curvature coupling is known for
electrons on a torus~\cite{encinosa2022}).

\textbf{[OBSERVATION]} None of these doors climbs past the third
nucleus.  Read down the atlas, geometry alone stops at the elliptic
rung.  The climbs come from somewhere else:\ from linking more densities
in one integral.

%% ========================================================================
\section{The climb:\ linking densities}
\label{sec:climb}
%% ========================================================================

\emph{Vocabulary.}  A two-center \emph{density} is the product of two
orbitals on two nuclei.  A \emph{kernel} joins two electrons:\ the
Coulomb kernel $1/r_{12}$, or a correlation factor (geminal)
$e^{-\gamma r_{12}}$.  An ordinary electron-repulsion integral links two
densities by one kernel.  Integrals of explicitly correlated methods link
more.  Paper~63 draws each integral as a \emph{density graph}:\ a vertex
for every electron and every nucleus, a line of mass $\zeta$ from a
nucleus to an electron for every orbital of exponent $\zeta$, and a line
between two electrons for every kernel, massless for the Coulomb kernel
and of mass $\gamma$ for a geminal~\cite{loutey_paper63}.  The
\emph{fibre} of an integral is the integral with all centers coincident;\
the class is read there, on a family in which one scale $t$ varies, from
the linear differential operator of least order that annihilates the
family.  The \emph{holomorphic period} is the period of the curve or
surface the densities build, the integral of the inverse dispersion roots
alone (the left side of Eq.~\eqref{eq:tree} below);\ the fibre of the
real integral carries further weights.  An integral sits on a
\emph{rung} when that operator has no irreducible piece (composition
factor) above the rung's.

\emph{The tree identity.}  \textbf{[SYMBOLIC + MEASURED]} Take $n$
two-center $1s$ densities joined into a tree by $n-1$ rational kernels.
In momentum space each density carries a dispersion root
$\Delta_v=\sqrt{c_vL_v^2+1}$, with $c_v$ the density's scale (Feynman)
parameter and $L_v$ the signed sum of the momenta on the kernels at $v$,
and Fourier-transforming each $1/\Delta_v$
turns every momentum integration into a delta function, so
\begin{equation}
  \int_{\mathbb{R}^{n-1}}\! d^{n-1}k\,\prod_{v=1}^{n}\frac{1}{\Delta_v}
  = \frac{2^{n}}{\pi\sqrt{c_1\cdots c_n}}
    \int_0^\infty \prod_{v=1}^{n} K_0(a_v x)\,dx ,
  \label{eq:tree}
\end{equation}
with $a_v=c_v^{-1/2}$, for a star and a chain alike:\ the holomorphic
period of a tree depends only on how many densities the kernels link, not
on the pattern of the links (Paper~63~\cite{loutey_paper63}).  The right
side is a Bessel moment in $n$ scales.

\emph{The operators.}  \textbf{[SYMBOLIC]} Vary one scale,
$F_n(t)=\int_0^\infty K_0(x)^{n-1}K_0(tx)\,dx$.  Then $F_n$ satisfies an
explicit linear equation in $t$ of order $n$, \textbf{[MEASURED]} each
checked against direct quadrature.  For $n=2$ it is the Legendre equation
of the complete elliptic integral $K$.  For $n=3$ it is the equation of
Clausen's square of an elliptic period.  For $n=4$, in $x=1/(4t^2)$ with
$\theta=x\,d/dx$, it is
\begin{equation}
  \theta^4 - x(2\theta+1)^2(10\theta^2+10\theta+3)
  + 9x^2(2\theta+1)^2(2\theta+3)^2 ,
  \label{eq:f4}
\end{equation}
\textbf{[PRIOR ART]} entry \#58 of the Calabi--Yau tables of Almkvist,
van Enckevort, van Straten and Zudilin~\cite{aesz2005}.
\textbf{[SYMBOLIC]} The operator of $F_4$ is irreducible, and $F_4$ is
not a polynomial in elliptic periods with algebraic moduli:\ two linked
densities give an elliptic period, three give a K3 surface whose period
is elliptic squared, and four give the first period of a new kind.  At
$n=5$ the tables print the operator verbatim, attached to entry \#39, and
its structure is that of the exterior square of a fourth-order
Calabi--Yau operator, so five linked densities give no new type of
period.  The operators of $F_n$ are irreducible for $n=2$ to $8$.

\emph{The banana integrals.}  \textbf{[SYMBOLIC]}
\gvc{p59-tree-banana-hadamard}{The tree ladder is tied to the banana
ladder, where an $l$-loop equal-mass banana integral carries a
Calabi--Yau $(l-1)$-fold~\cite{bonisch2021}, by one exact operation:\
with $y=1/(16t^2)$ and $F_n=t^{-1}g(y)$, the operator of $g$ is the
Hadamard product of the Legendre operator with the operator of the banana
with $n-1$ propagators, checked in exact arithmetic for $n=2$ to $5$}
(Paper~63~\cite{loutey_paper63}).  The same fact reads:\ $F_n$ is the
$n$-propagator banana, with masses $(t,1,\dots,1)$, integrated over a
line of external momentum.  It sits on the rung of the banana with
$n+1$ propagators but is not that family:\ its singular values of $t^2$,
$\{1\},\{4\},\{1,9\},\{4,16\}$ for $n=2,\dots,5$, are the thresholds of the
banana with $n-1$ propagators.  \textbf{[PRIOR ART]} The Hadamard
construction and the table entries are Almkvist and
Zudilin's~\cite{almkvist_zudilin2007,aesz2005};\ the operators for the
equal-mass two-dimensional banana integral are
Vanhove's~\cite{vanhove2014};\ the links of random-walk
moments to Bessel moments and Calabi--Yau equations are
Broadhurst's~\cite{broadhurst2009walks}.  At $t=1$ the ladder is the
Ising class of Bessel moments, $\int_0^\infty K_0(x)^n\,dx$, whose known values are prior art
(Bailey, Borwein, Broadhurst and Glasser~\cite{bbbg2008}).

\emph{Three linked densities.}  \textbf{[SYMBOLIC + MEASURED]} The chain
integral $\langle\rho_a(1)\rho_b(2)\rho_c(3)/(r_{12}r_{13})\rangle$, three
two-center $1s$ Slater densities joined by two Coulomb kernels at one
electron, has a branch locus of six lines with no four concurrent, so the
double cover is a K3 surface.  \textbf{[MEASURED + PRIOR ART]} Its
holomorphic period reduces to the three-scale Bessel moment
$\int_0^\infty K_0(at)K_0(bt)K_0(ct)\,dt$, whose closed form, a sum of
two products of complete elliptic integrals, is Eq.~33 of
Ref.~\cite{bbbg2008}, obtained there from an identity of
Bailey~\cite{bailey1936}.  The reduction from a Slater chain to that
moment is Paper~63's;\ the closed form is not.

\emph{What is proved about trees, and for which.}  Four statements
should be kept apart (all Paper~63~\cite{loutey_paper63}).

(a) \emph{The dictionary.}  \textbf{[SYMBOLIC PROOF + PRIOR ART]} Every
exponential-type orbital on a center is a polynomial differential
operator in its exponent and its center coordinate, applied to the $1s$
Yukawa function $e^{-\zeta s}/s$ on that center.  So every multi-density
integral with Coulomb or Slater-geminal kernels is that operator applied
to the $1s$ integral of the same density graph, and lies in the
$D$-module that integral generates (the functions reached from it by differentiating in its parameters and
multiplying by polynomial coefficients).  That alone does not bound the rung.

(b) \emph{The bound for the holomorphic period.}  \textbf{[SYMBOLIC
PROOF]} On the slice $(t;1,\dots,1)$ (one density's scale $t$ varies and
every other scale is one) of a tree of $n$ linked $s$-type densities, the holomorphic period is a differential operator applied to
$F_n$, so it sits at or below rung $n$;\ if $n\le8$ and the period is not
identically zero, it sits on rung $n$.

(c) \emph{The bound for the three-dimensional fibre.}  \textbf{[SYMBOLIC
PROOF]} The real three-dimensional integral of every tree of radial
densities reduces to one-dimensional integrals (for the Coulomb kernel
this is Newton's shell theorem).  For a Coulomb tree at the weights
$\rho_v(r)=a_v^3K_0(a_vr)/(2\pi^2)$, on the slice $(t;1,\dots,1)$, every
composition factor of the fibre's minimal operator has order at most
$n$, for every $n$;\ the same order bound holds for densities
differentiated at most twice in their center (and to any degree in their
scale).

(d) \emph{Attainment.}  \textbf{[SYMBOLIC PROOF]} For a Coulomb tree of
$n$ $s$-type densities at those weights, on the slice with $t$ on any one
vertex and every other scale equal to one, with $3\le n\le8$, $F_n$ lies
in the $D$-module of the fibre, so the minimal operator of the fibre has
a composition factor equivalent to the operator of $F_n$.  This holds for
every tree, not only for stars.  For $n\ge9$ the proof gives only that
the two modules meet in a nonzero submodule.  \textbf{[SYMBOLIC PROOF]}
The same holds, $3\le n\le8$, for densities that are homogeneous center
derivatives of degree at most two with a nonzero angular integral.
\textbf{[SYMBOLIC + MEASURED]} One step of the proof constructs an
operator with a controlled leading coefficient;\ that operator was
computed exactly for six rooted trees, modulo one prime for four more,
and modulo two primes for nine larger ones up to $n=8$.  The checks of
that step beyond $n=5$ are modulo primes, not exact.  For the
homogeneous-derivative case only axial densities were computed, nothing
with $\ell>0$ at $n\ge6$ in that step, and scale derivatives are not
covered.  Separately, attainment is computed for the chain of three and
for both trees of four, the star and the path, which sit on the rung of
four and skip the K3 rung (minimality bounded by a search modulo a prime,
not proved).

\emph{The rung rule and its caveats.}  \textbf{[OBSERVATION]} Read as a
Feynman graph, every instance fits one pattern:\ the class is set by the
shape of the density graph and the masses on its lines, and along a tree
it is set by the number of distinct two-center densities the kernels
couple, not by the number of nuclei or electrons.  Paper~63 states this
as a bounded conjecture with its falsifiers.  It is a pattern across the
tested instances, not a derivation, and it carries four caveats.

First, the rung names an operator, not a finite set of numbers.
\textbf{[OBSERVATION, evidence]} At equal scales the three-electron
fibre collapses to K3 periods and $\pi^3$, but at four generic scale
triples integer-relation searches at 340 and 400 digits find no relation
in five bases built from the K3 and elliptic periods, their derivatives
and products, dilogarithms, trilogarithms and weight-three Tate numbers,
while every basis recovers a planted relation (three of
the twenty searches are void);\ on both one-parameter
slices computed the fibre is a K3 function integrated once.  That is
evidence that the generic fibre carries content beyond the K3 and
elliptic periods, not a theorem.

Second, rings are ordered by uniqueness, not by count.
\textbf{[SYMBOLIC + MEASURED]} A Coulomb ring of three densities, the
triangle $\langle\rho_a\rho_b\rho_c/(r_{12}r_{13}r_{23})\rangle$, does
not climb:\ its loop integrates out exactly by the three-dimensional
uniqueness (star-triangle) relation, and it sits on the elliptic rung,
below the chain of three.  At one center it reproduces Remiddi's closed
form~\cite{remiddi1991,ppr2004}, which carries a dilogarithm where the
chain is rational.  The correlated triangle of explicitly correlated
methods, with exponential geminals on two edges, does not drop:\ it sits
with the chain on the K3 rung at $\gamma=1$, $3/2$ and $2$, the geminal
acting as an imaginary external momentum.  A Gaussian geminal leaves the
ladder for exponential periods.  Only the triangle was computed;\ the
ring of four stopped for feasibility.

Third, the rung is read on the fibre and its holomorphic period;\ the
integrated value, over the separations and Feynman parameters, may sit
lower, and in one class it does:\ the two-center exchange integral builds the elliptic fibre
and closes at genus zero.

Fourth, the statements above hold where they are proved:\ for $s$-type
densities, on one-scale slices, with the weights named.  That every
computed general-orbital fibre sits on its graph's $1s$ rung is observed,
not proved.

\emph{What this means for a molecule.}  \textbf{[OBSERVATION]}
\gvc{p63-hamiltonian-elliptic-stype}{Every electron-repulsion integral
couples exactly two densities, so over $s$-type orbitals the integrals of
a molecular Hamiltonian stay on the elliptic rung however large the
molecule (for $\ell>0$ this is not shown)}.  The class climbs only in
correlated integrals that link more densities
(Paper~63~\cite{loutey_paper63}).

%% ========================================================================
\section{Two centers in detail}
\label{sec:two}
%% ========================================================================

The first door of the chemistry is the best understood, and most of it
is classical.  Paper~58 evaluates all four classes of two-center
electron-repulsion integral in closed form and records what each closed
form contains~\cite{loutey_paper58}:

\begin{center}
\small
\begin{tabular}{@{}ll@{}}
\toprule
class & transcendental content \\
\midrule
one-center, $(AA|BB)$ & none:\ elementary \\
$(AA|AB)$, $(AB|BB)$ & $\{E_1,\mathrm{Ei},\ln\}$;\ $\gamma$ at equal rates;\\
 & none if the $AA$ pair is $s$-type \\
$(AB|AB)$ & $\{E_1,\ln,\gamma\}$ per Neumann term \\
\bottomrule
\end{tabular}
\end{center}

\noindent \textbf{[PRIOR ART]} The first and last rows are known:\ the
two-center Coulomb integral over Slater orbitals reduces to elementary
functions, and the exchange integral's closed form carries $E_1$, $\ln$
and Euler's $\gamma$ (Lesiuk and
Moszynski~\cite{lesiuk2014a,lesiuk2014b}, who attribute the general
exchange form to Maslen and Trefry;\ for $1s$, Sugiura and
Ruedenberg~\cite{sugiura1927,ruedenberg1951}).  Those references evaluate
the hybrid class by one-dimensional quadrature;\ Paper~58 did not find
its closed form stated.  No $\pi$ survives in any of them:\ it enters
through the spherical normalizations and cancels against the $4\pi$ of
the Coulomb kernel (Paper~18~\cite{loutey_paper18}).

\emph{Weight.}  The exchange integral is computed as a series, the
Neumann expansion.  Each term is an iterated integral over a simplex, the
shape that would generically give weight two, and it does not:\
assembled in closed form, and checked against quadrature, it contains
only $\exp$, $E_1$, $\ln$ and $\gamma$, so the seed sets form a
weight-one filtration for every Neumann term, and for the whole integral
at equal exponents, where the series stops.  At unequal exponents the series does not stop,
and the sum carries the two-center master integral $f_0$, which Pachucki
obtained from a second-order differential equation in the internuclear
distance~\cite{pachucki2009}.  \textbf{[SYMBOLIC PROOF]} For $0<a<b$,
$f_0$ is a two-fold iterated integral with an explicit closed form;\ its
dilogarithm constants cancel, because $f_0(0)=0$;\ and it is not a
combination, with coefficients rational in $R$, of exponentials,
logarithms and exponential integrals for any finite set of rates.  As
functions of $R$, the exchange integrals whose order nine is proved
below therefore have iterated-integral weight two, not one, for $0<a<b$.
At $a=b$ the weight-two part degenerates to weight one, consistent with
Sugiura's formula.

\emph{Exact orders.}  \textbf{[SYMBOLIC]} Each class obeys a linear
differential equation in $R$ with integer polynomial coefficients;\ the
transcendentals live only in its solutions.  \textbf{[SYMBOLIC PROOF]}
For an integral that contains only exponentials (the one-electron
integrals and $(AA|BB)$), the order of the least operator is the number
of distinct decay rates with a nonzero coefficient, at every choice of
exponents.  For exchange the order is exactly $4$ at equal exponents and
exactly $9$ for every $a\neq b$, proved for seven integrals:\ $1s$, four
integrals with $p$ or $d$ orbitals and $\sigma=0$, and the two $m\neq0$
integrals against $1s$ (those two given Pachucki's first-order equation in the
$r_{12}$ exponent, used as published).  The bound comes from differentiating Pachucki's
equation at the exchange point (for the two $m\neq0$ integrals, from his
masters weighted by powers of $r_{12}$), which places each integral in the span
of nine explicit functions closed under $d/dR$;\ exactness comes from the
independence of those nine functions and an exact determinant test.  For
hybrids with a multipole $L\ge1$ on the same-center pair, that the order found is also the least is proved at every $Z_A,Z_B>0$
for the nine hybrids of the panel with an $L\ge1$ pair and for
$(2p_0 2p_0|1s\,1s_{\rm B})$:\ the order is $5$ away from the two
coincidences of decay rates and $4$ at each, with no other exceptional
ratio.  This rests on the independence of the atoms (a hand proof,
reviewed once), on the counting of the order (a hand proof, reviewed
once) and on the parametric closed forms, which agree with the evaluator
at sample points in every region;\ for four of the ten integrals the
analysis was reviewed (two of them also against an independent
quadrature), while for the other six the same pipeline was run without a
separate review.

\emph{What the closed forms decide.}
Because the direct Coulomb block is elementary, its entries are not
merely counted but decided:\ all $195$ permitted entries of that block are proven nonzero by Lindemann--Weierstrass separation.
\textbf{[SYMBOLIC + MEASURED]} And the
potential curve of minimal-basis H$_2$ is a single closed-form expression
in $\{\exp,E_1,\log,\gamma\}$.

\emph{What is open at two centers.}  Whether $f_0$ has a dilogarithm or
product form, and what holds for its values at a fixed $R$, is not
decided.  The weight-two argument is given for $0<a<b$;\ the same
argument on the real branch should cover $a>b$, but that is not checked.
And Paper~58 states that its closed form for $f_0$ was not compared with
the classical unequal-exponent literature beyond
Refs.~\cite{pachucki2009,lesiuk2014b}, so no novelty is claimed for it.

%% ========================================================================
\section{Three centers, and water}
\label{sec:three}
%% ========================================================================

Prolate spheroidal coordinates admit exactly two foci.  Paper~59
evaluates the two-electron three-center integral $(XY|XZ)$ over $1s$
orbitals in momentum space, where each center enters as a phase rather
than a focus~\cite{loutey_paper59}.  There the two densities carry
dispersion roots with different scales, and their product is the quartic
\begin{equation}
  y^{2}=(c_{1}k^{2}+1)(c_{2}k^{2}+1),
  \label{eq:curve}
\end{equation}
an elliptic curve.  \textbf{[MEASURED / classical]} The radial content
of the integral is a two-scale Bessel moment whose period at zero
separation is a complete elliptic integral,
\begin{equation}
  \int_{0}^{\infty}\!\frac{dk}{\sqrt{(c_{1}k^{2}+1)(c_{2}k^{2}+1)}}
  =\frac{K(m)}{\sqrt{c_{\max}}},\quad m=1-\frac{c_{\min}}{c_{\max}},
  \label{eq:period}
\end{equation}
verified to 31 digits, degenerating to the elementary two-center case
exactly on the diagonal $c_1=c_2$.  This is the case $n=2$ of
Eq.~\eqref{eq:tree}.  \textbf{[SYMBOLIC + MEASURED]} The one-mass slice
satisfies an explicit fourth-order Picard--Fuchs equation (the linear
differential equation a family of periods obeys), and its
connection is irreducible, so no closed form arises through
factorization.  \textbf{[OPEN]} Paper~59 does not give a finite closed
form, and shows that the natural route to one, the elliptic-dilogarithm
mechanism of the unequal-mass sunrise~\cite{blochvanhove2015}, does not
engage.  The message of that paper is a relocation, not a solution:\ the
third center raises the radial core's class by one genus.

\emph{The number theory of the radial core.}  Paper~61 develops
it~\cite{loutey_paper61}.  \textbf{[MEASURED;\ classical]} The elliptic
family is the Legendre family over the modular curve $X(2)$, and the
parameter integral passes through complex-multiplication fibres where
the periods are Chowla--Selberg
$\Gamma$-values~\cite{chowla_selberg1967}.  \textbf{[SYMBOLIC +
MEASURED]} The connection has four master periods, and its monodromy
group lies in $\mathrm{Sp}_4(\mathbb{Z})$.  The paper's guarded searches
are negative for a clean low-height closed form through weight three,
and the class of the integrated value remains its open frontier.

\emph{Water, and the bond angle.}  The wall is closed-form, not
computational.  \textbf{[MEASURED]} On the water geometry, with Slater
orbitals, seven three-center integrals of both types converge to about
$10^{-9}$ by $L\approx12$ in the multipole series about the shared
center (the Barnett--Coulson route~\cite{barnett_coulson1951}).  That
series is the Legendre series of the integral in the cosine $x$ of the
bond angle, so the decay of its terms locates the nearest singular point
of the integral as a function of the angle (all
Paper~59~\cite{loutey_paper59}):

\begin{itemize}
\item \textbf{[MEASURED]} At water's equal bond lengths the terms fall
like $\ell^{-8}$, which puts the singular point where the two hydrogens
would merge.  Bent water is an ordinary point.
\item \textbf{[SYMBOLIC + MEASURED]} The singular term is derived.  Only
the two hydrogen cusps, meeting through $1/r_{12}$, make the integral
non-analytic in the angle, and the lowest such term is $C\,|Y-Z|^{7}$
with $C=\zeta_{\rm H}^{5}\zeta_{\rm X}^{3}e^{-\zeta_{\rm
X}(D_1+D_2)}/630$.  The two-electron term was found first by
Matuzas~\cite{matuzas2026sc}, whose preprint Paper~59 found only after
its own work.
\item \textbf{[SYMBOLIC + MEASURED]} The bond-angle function obeys a
linear differential equation in $x$.  For the one-electron three-center
integral it has order $8$ with polynomial coefficients of degree $31$,
and its only finite singular point with a non-integer exponent is the
point $x^{*}$ where the two outer centers merge, with exponents $3/2$ and
$5/2$.  For $(XY|XZ)$ every Legendre coefficient is a rational
combination of 31 fixed constants, each an exponential times $1$, $E_1$
or a logarithm, and the coefficients obey a recurrence in $\ell$ of order
$41$ (one geometry, two primes);\ no differential operator was found
within the size searched.  The constants are assumed linearly
independent over $\mathbb{Q}$.  For
the other three-center type, a one-center density against a two-center
one, the fixed set has $21$ constants and the minimal operator found (not
certified) has order $14$, with non-integer exponents $3/2$, $7/2$ and
$9/2$ at $x^{*}$.
\item \textbf{[MEASURED]} The integrated three-center value leaves the
two-center class as a function of the bond angle:\ its bond-angle
operator keeps a piece of order six with no first-order factor at two
geometries, where a two-center-class form would factor completely
(assuming its constants independent;\ the class of the value at one
geometry is not decided).
\end{itemize}

\noindent These statements are for $1s$ functions at the geometries and
exponents each record names.

%% ========================================================================
\section{What this is not}
\label{sec:not}
%% ========================================================================

\emph{It is not faster chemistry.}  Knowing the class of an integral
does not make a molecular energy more accurate.  Paper~59 says so of its
own result:\ it does not by itself deliver a polyatomic energy, and what
the reduction buys is an evaluator and a diagnosis of why the integral
is hard~\cite{loutey_paper59}.  The series' qubit-resource comparisons
(Paper~14) are at matched qubit counts, not matched
accuracy~\cite{loutey_paper14}.
The one matched-accuracy comparison is Paper~62's, on water with tuned exponential bases, and it depends on the Gaussian family:\
\textbf{[MEASURED, pre-registered]}
\gvc{p62-matched-accuracy-ceiling}{at matched accuracy the Gaussian cost
depends on the family.  In the cc-pV family aug-cc-pVDZ at 41 functions
stays $\gvq{p62_aug_above_dzp_mha}{8.0}\mHa$ short of the Slater DZP
while already costing $\gvq{p62_aug_pauli_ratio}{9.2}\times$ the Pauli
terms.  In the ANO family a 35-function truncation of ANO-R2 is
$\gvq{p62_a35_below_dzp_mha}{6.2}\mHa$ below it, at
$\gvq{p62_a35_pauli_ratio}{4.7}\times$ the Pauli terms.  The measured
matched-accuracy cost is therefore at most $4.7\times$ the terms}
(Paper~62~\cite{loutey_paper62}).  Hydrogen fluoride, a second molecule,
repeats the cc-pV finding;\ the ANO family was not run there.  A ring of
six hydrogens, with no core, does not, so
\gvc{p62-hring-contrast}{the cc-pV matched-accuracy margin is a property of these two
molecules, not of exponential functions}.  That paper adds that nothing in it is
a lattice structure of the series:\ its bases are tuned exponential
functions evaluated through Gaussian fits.  Paper~18 reads the elliptic
curve of Sec.~\ref{sec:three} as the reason Gaussian orbitals, whose
product theorem collapses three centers to one, displaced Slater orbitals
for polyatomic quantum chemistry~\cite{loutey_paper18}.

\emph{It is not a new foundation.}  The sphere is Fock's.  The
equivalence between the label set with its Casimir and the
Schr\"odinger problem is carried by his continuum projection, and
Paper~7 audits that chain rather than replacing
it~\cite{loutey_paper7}.  The words ``forced'' and ``price'' are
catalogue vocabulary.  And the ladder itself is not new mathematics:\
the operators on it are known (Sec.~\ref{sec:prior}).

%% ========================================================================
\section{Prior art, and what the series adds}
\label{sec:prior}
%% ========================================================================

\emph{The sphere and the rule.}  Fock's projection~\cite{fock1935};\
the Peter--Weyl theorem~\cite{peter_weyl1927};\ the compactness reading
of Liu and Noui~\cite{liu_noui2017} and of
Rovelli~\cite{rovelli1996};\ the failure of its converse, Maz'ya and
Shubin~\cite{mazya_shubin2005};\ periods as a class of numbers,
Kontsevich and Zagier~\cite{kontsevich_zagier2001}.

\emph{Two centers.}  Sugiura's $1s$ exchange
integral~\cite{sugiura1927};\ Roothaan's and Ruedenberg's two-center
integrals~\cite{roothaan1951,ruedenberg1951};\ the general exchange
form, attributed by Lesiuk and Moszynski to Maslen and Trefry, and
Lesiuk and Moszynski's own treatment of the Coulomb, hybrid and exchange
classes~\cite{lesiuk2014a,lesiuk2014b};\ Pachucki's differential
equation for the two-center master integral~\cite{pachucki2009}.

\emph{Three centers.}  Shibuya and Wulfman's momentum-space molecular
orbitals~\cite{shibuya1965};\ the Barnett--Coulson
expansion~\cite{barnett_coulson1951};\ the Averys' momentum-space
Coulomb--Sturmian program, which treats exactly these integrals under
the same Fock projection and reaches them through an infinite series in
a one-center basis~\cite{avery_avery_book2018,avery2013};\ the elliptic
dilogarithm of the sunset graph, Bloch and
Vanhove~\cite{blochvanhove2015}, and the sunrise as a higher normal
function, Bloch, Kerr and Vanhove~\cite{blochkerrvanhove2015};\ the
bond-angle cusp term, Matuzas~\cite{matuzas2026sc}.

\emph{The ladder.}  The two-scale and three-scale Bessel moments and
their elliptic evaluations, Bailey, Borwein, Broadhurst and
Glasser~\cite{bbbg2008}, on an identity of Bailey~\cite{bailey1936};\
random-walk moments and Calabi--Yau equations,
Broadhurst~\cite{broadhurst2009walks};\ the Hadamard construction and
the Calabi--Yau tables, Almkvist and Zudilin and their
coauthors~\cite{almkvist_zudilin2007,aesz2005};\ the operators of the
equal-mass two-dimensional banana integral,
Vanhove~\cite{vanhove2014};\ the Calabi--Yau $(l-1)$-fold of
the $l$-loop banana, B\"onisch, Fischbach, Klemm, Nega and
Safari~\cite{bonisch2021};\ the Kloosterman motives behind the Ising
class, Fres\'an, Sabbah and Yu~\cite{fresansabbahyu2023};\ the atomic
three-electron integral, Remiddi~\cite{remiddi1991}.
The nearest theorem in shape to the rung rule is Brown's:\ for massless
primitive $\phi^4$ graphs in four dimensions the set of periods is minor
monotone~\cite{brown2009periods}, with an extension to generic masses in
even dimension that is a conjecture with an affine form
proved~\cite{brown2015coaction}.  The lines here carry special masses in
three dimensions, outside both hypotheses.  Matuzas~\cite{matuzas2026} relates one channel of a different
four-electron correlation integral, a ring of Coulomb links with one
diagonal over Slater functions, to another Calabi--Yau operator,
AESZ-330, by an inhomogeneous equation.

\emph{What the series adds,} as its papers state it:
\begin{itemize}
\item the accounting:\ each released axis carries a named number class
(Paper~18), kept as one register with a test behind each row;
\item at two centers:\ the use of the closed-form classes to decide
which integrals vanish;\ the least operator in $R$ of each class, with its order
proved (exactly $9$ and exactly $4$ for exchange) and its singular points
read, for $1s$ orbitals plus the $l\le2$ panel;\ and the proof that the
seven exchange integrals whose order $9$ is proved have weight two as
functions of $R$ for $0<a<b$ (Paper~58;\ novelty against the classical
literature is not claimed for the last);
\item at three centers:\ the reading of the radial core as an elliptic
Bessel moment, a connection a targeted literature search found unmade in
either field \textbf{[OBSERVATION]};\ the fixed constant sets of the
bond-angle series, and its differential equations, which were not found
in a literature search (Paper~59);
\item on the ladder:\ the route from molecular integrals to the known
operators, the tree identity~\eqref{eq:tree}, the dictionary that places every orbital's integral in the $D$-module of
its graph's $1s$ integral, the readings of the ring, the geminal and the
Gaussian, and the fact that most rungs were registered in committed
records before they were computed (Paper~63).  Paper~63 also proves that
trees of up to eight $s$-type densities reach their rung, on the slices
and at the weights of Sec.~\ref{sec:climb}.  That no source places a molecular integral on
these operators is an absence from a targeted search, not a proof of
absence.
\end{itemize}

%% ========================================================================
\section{Open questions}
\label{sec:open}
%% ========================================================================

These are handed over as the owning papers leave them.

\begin{enumerate}
\item \emph{Trees with other densities.}  Whether the three-dimensional
fibre reaches rung $n$ for trees whose densities are non-homogeneous
center derivatives, scale derivatives, or have a vanishing angular
integral.  Paper~63 states the first as a conjecture and gives a local
argument for it, not a proof.
\item \emph{Nine or more densities.}  For $n\ge9$ the attainment proof
gives only a nonzero common submodule, because the irreducibility of the
operator of $F_n$ is established for $n\le8$.
\item \emph{The ring of four.}  Among rings only the triangle was
computed.  Whether each Coulomb loop that integrates out by uniqueness
lowers the rung by one density, as the conjecture says, has one instance
behind it.
\item \emph{The two-center master integral.}  Its weight for $a>b$;\
whether it has a dilogarithm or product form;\ the class of its values
at a fixed $R$ (Paper~58).
\item \emph{The generic three-electron fibre.}  Whether a principled
widening of the basis finds a relation where the present searches find
none.  A relation would show that the fibre sits in a finite ring of
periods after all;\ Paper~63 names this as the case that tests its
conjecture.
\item \emph{The first rung's integrated value.}  A finite closed form
for the three-center integral, or a proof that none exists in a named
class (Papers~59 and~61).
\item \emph{The high-field Zeeman class.}  The number class of a
hydrogen eigenvalue at finite magnetic field (atlas, door 29).
\item \emph{A proof of the rule.}  Paper~63 expects that a proof would
run through the Symanzik polynomials of the density graph with the
Yukawa masses on its lines, and has not attempted it.
\end{enumerate}

%% ========================================================================
\section{How to check any of this}
\label{sec:check}
%% ========================================================================

Most claims above have a test in the repository's \texttt{tests/}
directory, and \texttt{docs/claims\_register.md} lists the headline
claims with their tiers;\ the three-electron equal-scale collapse and
generic null are carried by prediction records in
\texttt{docs/predictions/}, not tests.  The skeleton and its wall:\
\texttt{test\_fock\_projection.py}, \texttt{test\_fock\_laplacian.py},
\texttt{test\_paper7\_label\_set.py},
\texttt{test\_graph\_s3\_nonconvergence.py}.  Two centers:\
\texttt{test\_paper58\_abelian\_residue.py},
\texttt{test\_paper58\_exact\_order.py},
\texttt{test\_paper58\_minimal\_order.py},
\texttt{test\_paper58\_master\_weight.py},
\texttt{test\_paper58\_exchange\_reverse.py},
\texttt{test\_paper58\_hybrid\_every\_ratio.py}.  Three centers:\
\texttt{test\_routeC\_momentum.py},
\texttt{test\_paper59\_water\_bond\_angle.py},
\texttt{test\_paper59\_aa\_bc\_minimal\_operator.py}.  The ladder:\
\texttt{test\_paper63\_dictionary.py},
\texttt{test\_paper63\_tree\_bound.py},
\texttt{test\_paper63\_tree\_reduction.py},
\texttt{test\_paper63\_tree\_rung.py},
\texttt{test\_paper63\_four\_density\_rung.py},
\texttt{test\_paper63\_three\_electron\_k3.py},
\texttt{test\_paper63\_four\_centre\_genus.py},
\texttt{test\_paper63\_tree\_attainment.py},
\texttt{test\_paper63\_tree\_attainment\_big.py},
\texttt{test\_paper63\_tree\_attainment\_l.py},
\texttt{test\_paper63\_density\_ring.py},
\texttt{test\_paper63\_gy\_gaussian.py}.  The doors:\
\texttt{test\_atlas\_register.py},
\texttt{test\_atlas\_gap\_doors.py},
\texttt{test\_renner\_teller\_m\_splitting.py},
\texttt{test\_paper62\_hydrogen\_torus.py},
\texttt{test\_paper62\_squashed\_ring.py}.  The matched-accuracy
comparison:\ \texttt{test\_paper62\_matched\_accuracy.py},
\texttt{test\_paper62\_ano\_matched\_accuracy.py}.  The predictions that
were written down before their computations are in
\texttt{docs/predictions/}, each with its outcome, including the ones
that missed.

%% ========================================================================
\begin{thebibliography}{99}

\bibitem{loutey_paper0}
J.~Loutey, ``Counting Angular Momentum Labels by Isotropic Packing,''
GeoVac Paper~0 (2026).

\bibitem{loutey_paper7}
J.~Loutey, ``The Dimensionless Vacuum:\ Recovering the Schr\"odinger
Equation on a Scale-Invariant Label Lattice,'' GeoVac Paper~7 (2026).

\bibitem{loutey_paper14}
J.~Loutey, ``Structurally Sparse Qubit Hamiltonians from Natural-Geometry
Hydrogenic Bases,'' GeoVac Paper~14 (2026).

\bibitem{loutey_paper18}
J.~Loutey, ``Spectral-Geometric Exchange Constants:\ Projecting Discrete
Spectra onto Continuous Manifolds,'' GeoVac Paper~18 (2026).

\bibitem{loutey_paper28}
J.~Loutey, ``Periods of Dirac Spectral Sums on $S^3$,'' GeoVac Paper~28
(2026).

\bibitem{loutey_paper32}
J.~Loutey, ``The GeoVac Spectral Triple:\ A Synthesis Paper Locking the
Framework into the Marcolli--van~Suijlekom Gauge-Network Lineage,''
GeoVac Paper~32 (2026).

\bibitem{loutey_paper34}
J.~Loutey, ``GeoVac as a Two-Layer Framework:\ Projection Taxonomy and
the Empirical Matches Catalog,'' GeoVac Paper~34 (2026).

\bibitem{loutey_paper55}
J.~Loutey, ``Periods of GeoVac:\ a cyclotomic mixed-Tate census of the
master Mellin engine,'' GeoVac Paper~55 (2026).

\bibitem{loutey_paper58}
J.~Loutey, ``Angular Sparsity Is an Atomic-Sector Property:\ The Abelian
Residue of Multi-Center Bonding,'' GeoVac Paper~58 (2026).

\bibitem{loutey_paper59}
J.~Loutey, ``The Three-Center Electron-Repulsion Integral Is an Elliptic
Bessel Moment:\ A Genus-One Radial Core at the Third Center in
Momentum-Space Slater Theory,'' GeoVac Paper~59 (2026).

\bibitem{loutey_paper61}
J.~Loutey, ``The Bessel-Moment Period Algebra of the Three-Center Coulomb
Integral:\ Modular Structure and Cosmic-Galois Placement,'' GeoVac
Paper~61 (2026).

\bibitem{loutey_paper62}
J.~Loutey, ``Water at Twenty-Four Functions:\ Slater-Type Bases Set on
the Free Atom Against cc-pVDZ, With the Qubit Cost,'' GeoVac Paper~62
(2026).

\bibitem{loutey_paper63}
J.~Loutey, ``Linked Electron Densities Climb the Banana Ladder:\ Period
Classes of Molecular Integrals over Exponential-Type Orbitals,'' GeoVac
Paper~63 (2026).

\bibitem{fock1935}
V.~Fock, ``Zur Theorie des Wasserstoffatoms,''
\textit{Z.\ Phys.}\ \textbf{98}, 145--154 (1935).

\bibitem{peter_weyl1927}
F.~Peter and H.~Weyl,
``Die Vollst\"andigkeit der primitiven Darstellungen einer
geschlossenen kontinuierlichen Gruppe,''
\textit{Math.\ Ann.}\ \textbf{97}, 737--755 (1927).

\bibitem{liu_noui2017}
C.~Liu and K.~Noui, ``Gravity as an $SU(1,1)$ gauge theory in four
dimensions,'' \textit{Class.\ Quantum Grav.}\ \textbf{34}, 135008 (2017),
arXiv:1702.06793.

\bibitem{rovelli1996}
C.~Rovelli, ``Relational quantum mechanics,''
\textit{Int.\ J.\ Theor.\ Phys.}\ \textbf{35}, 1637--1678 (1996),
arXiv:quant-ph/9609002.

\bibitem{mazya_shubin2005}
V.~Maz'ya and M.~Shubin, ``Discreteness of spectrum and positivity criteria
for Schr\"odinger operators,''
\textit{Ann.\ of Math.}\ \textbf{162}, 919--942 (2005).

\bibitem{kontsevich_zagier2001}
M.~Kontsevich and D.~Zagier, ``Periods,'' in \textit{Mathematics
Unlimited --- 2001 and Beyond}, B.~Engquist and W.~Schmid, eds.,
Springer (2001), pp.~771--808.

\bibitem{renner1934}
R.~Renner, ``Zur Theorie der Wechselwirkung zwischen Elektronen- und
Kernbewegung bei dreiatomigen, stabf\"ormigen Molek\"ulen,''
\textit{Z.~Phys.}\ \textbf{92}, 172--193 (1934).

\bibitem{encinosa2022}
M.~Encinosa and J.~Williamson, ``Electron gases in toroidal shells: mode
coupling and state functions,'' Physica E \textbf{146}, 115558 (2023),
doi:10.1016/j.physe.2022.115558; arXiv:2207.01745.

\bibitem{sugiura1927}
Y.~Sugiura, ``\"Uber die Eigenschaften des Wasserstoffmolek\"uls im
Grundzustande,'' Z.~Phys.\ \textbf{45}, 484--492 (1927).

\bibitem{roothaan1951}
C.~C.~J.~Roothaan, ``A study of two-center integrals useful in
calculations on molecular structure.\ I,''
\textit{J.\ Chem.\ Phys.}\ \textbf{19}, 1445--1458 (1951).

\bibitem{ruedenberg1951}
K.~Ruedenberg, ``A study of two-center integrals useful in calculations
on molecular structure.\ II.\ The two-center exchange integrals,''
\textit{J.\ Chem.\ Phys.}\ \textbf{19}, 1459--1477 (1951).

\bibitem{pachucki2009}
K.~Pachucki, ``Two-center two-electron integrals with exponential functions,''
\textit{Phys.\ Rev.\ A}\ \textbf{80}, 032520 (2009).

\bibitem{lesiuk2014a}
M.~Lesiuk and R.~Moszynski, ``Calculation of two-centre two-electron integrals
over Slater-type orbitals revisited.\ I.\ Coulomb and hybrid integrals,''
\textit{Phys.\ Rev.\ E} \textbf{90}, 063318 (2014), arXiv:1410.6112.

\bibitem{lesiuk2014b}
M.~Lesiuk and R.~Moszynski, ``Calculation of two-centre two-electron integrals
over Slater-type orbitals revisited.\ II.\ Neumann expansion of the exchange
integrals,'' \textit{Phys.\ Rev.\ E} \textbf{90}, 063319 (2014), arXiv:1410.6113.

\bibitem{shibuya1965}
T.~Shibuya and C.~E.~Wulfman, ``Molecular orbitals in momentum space,''
\textit{Proc.\ R.\ Soc.\ Lond.\ A}\ \textbf{286}, 376 (1965).

\bibitem{barnett_coulson1951}
M.~P.~Barnett and C.~A.~Coulson, ``The evaluation of integrals occurring
in the theory of molecular structure.\ Parts I \& II,''
\textit{Philos.\ Trans.\ R.\ Soc.\ Lond.\ A}\ \textbf{243}, 221--249 (1951).

\bibitem{avery_avery_book2018}
J.~E.~Avery and J.~S.~Avery, \textit{Hyperspherical Harmonics and Their
Physical Applications} (World Scientific, Singapore, 2018).

\bibitem{avery2013}
J.~E.~Avery, ``Fast electron repulsion integrals for molecular Coulomb
Sturmians,'' \textit{Adv.\ Quantum Chem.}\ \textbf{67}, 129--151 (2013).

\bibitem{matuzas2026sc}
J.~Matuzas, ``Single-centre expansions of three-centre 1s Slater integrals:\ convergence rates
and certified error bounds,'' Zenodo preprint, doi:10.5281/zenodo.23448184 (2026).

\bibitem{chowla_selberg1967}
A.~Selberg and S.~Chowla, ``On Epstein's zeta-function,''
\textit{J.\ Reine Angew.\ Math.}\ \textbf{227}, 86 (1967).

\bibitem{blochvanhove2015}
S.~Bloch and P.~Vanhove, ``The elliptic dilogarithm for the sunset graph,''
\textit{J.\ Number Theory}\ \textbf{148}, 328--364 (2015); arXiv:1309.5865.

\bibitem{blochkerrvanhove2015}
S.~Bloch, M.~Kerr, and P.~Vanhove, ``A Feynman integral via higher normal
functions,'' \textit{Compos.\ Math.}\ \textbf{151}, 2329--2375 (2015);
arXiv:1406.2664.

\bibitem{bbbg2008}
D.~H.~Bailey, J.~M.~Borwein, D.~Broadhurst, and M.~L.~Glasser, ``Elliptic
integral evaluations of Bessel moments and applications,''
\textit{J.\ Phys.\ A}\ \textbf{41}, 205203 (2008); arXiv:0801.0891.

\bibitem{bailey1936}
W.~N.~Bailey, ``Some infinite integrals involving Bessel functions (II),''
J.\ London Math.\ Soc.\ \textbf{11}, 16--20 (1936).

\bibitem{broadhurst2009walks}
D.~Broadhurst, ``Bessel moments, random walks and Calabi--Yau equations,''
preprint (2009).

\bibitem{aesz2005}
G.~Almkvist, C.~van Enckevort, D.~van Straten, and W.~Zudilin, ``Tables of
Calabi--Yau equations,'' arXiv:math/0507430 (2005).

\bibitem{almkvist_zudilin2007}
G.~Almkvist and W.~Zudilin, ``Differential equations, mirror maps and zeta
values,'' in \textit{Mirror Symmetry V}, AMS/IP Studies in Advanced
Mathematics \textbf{38}, 481--515 (2007); arXiv:math/0402386.

\bibitem{vanhove2014}
P.~Vanhove, ``The physics and the mixed Hodge structure of Feynman
integrals,'' arXiv:1401.6438 (2014).

\bibitem{bonisch2021}
K.~B\"onisch, F.~Fischbach, A.~Klemm, C.~Nega, and R.~Safari, ``Analytic
structure of all loop banana integrals,'' \textit{J.\ High Energy Phys.}\
\textbf{05}, 066 (2021); arXiv:2008.10574.

\bibitem{fresansabbahyu2023}
J.~Fres\'an, C.~Sabbah, and J.-D.~Yu, ``Quadratic relations between Bessel
moments,'' \textit{Algebra Number Theory}\ \textbf{17}, 541--602 (2023);
arXiv:2006.02702.

\bibitem{brown2009periods}
F.~C.~S.~Brown, ``On the periods of some Feynman integrals,'' arXiv:0910.0114 (2009).

\bibitem{brown2015coaction}
F.~Brown, ``Feynman amplitudes, coaction principle, and cosmic Galois
group,'' arXiv:1512.06409 (2015).

\bibitem{remiddi1991}
E.~Remiddi, ``Analytic value of the atomic three-electron correlation
integral with Slater wave functions,'' \textit{Phys.\ Rev.\ A}\ \textbf{44},
5492 (1991).

\bibitem{ppr2004}
K.~Pachucki, M.~Puchalski and E.~Remiddi, ``Recursion relations for
Hylleraas three-electron integral,'' \textit{Phys.\ Rev.\ A} \textbf{70},
032502 (2004);\ arXiv:physics/0405057.

\bibitem{matuzas2026}
J.~Matuzas, ``A Calabi--Yau period in a four-electron correlation integral:
modular physical fibre and a hypergeometric endpoint,'' Zenodo,
doi:10.5281/zenodo.23218203 (2026).

\end{thebibliography}

\end{document}
